\section{Combined interpretation}
\label{sec:interpretation}

\subsection{What the direct conversion now means}

The successful object is not a single equation that turns a halo history into
one deterministic galaxy profile. It is the composition
\begin{equation}
  \Mpeak(t),\ctwo(t)
  \longrightarrow \vect{x}(z)
  \longrightarrow P(\vect{\theta}_{\cog}\mid\vect{x},z)
  \longrightarrow P[M_\star(<R)\mid\vect{x},z].
  \label{eq:final_composition}
\end{equation}
The first arrow defines a compact halo state, the second is a fitted conditional
distribution, and the third is a monotone coordinate decoder. The model is
interpretable because its intermediate galaxy variables are total mass,
central fraction, and radial quantiles. It remains statistical because much of
the profile-shape variance is unresolved by the selected halo inputs.

\subsection{Why readable coordinates can match PCA}

PCA is optimized to reconstruct variance in the particular training sample.
The readable coordinates are instead chosen to preserve scientifically useful
functionals of a monotone cumulative profile. Their success indicates that, on
the adopted radial grid, the dominant profile variation can be captured by one
normalization, one central anchor, and three ordered outer scales. The result
does not imply that these are unique sufficient statistics. Adding $R_{90}$
improves the representation floor but not the halo prediction, demonstrating
that sufficiency for profile reconstruction and predictability from halos are
different requirements.

\subsection{Why the model must remain probabilistic}

The conditional mean predicts total stellar mass well but explains only
15--36 per cent of the variance in individual shape coordinates at $z=0.4$.
This remaining scatter is physically plausible: baryonic feedback, merger
orbits, satellite structure, stellar stripping, and projection need not be
fully determined by a smooth main-branch MAH and one concentration measurement.
The correlated residual distribution restores most of the population width
that the mean suppresses. Removing it would change the scientific model, not
merely its error bars.

\subsection{What the epoch-local result changes}

The strongest \ExpFiftyOne{} result is not simply that a $z=2$ fit succeeds. It
is that the fit improves when the halo assembly parameters are defined using
the correct epoch-local clock. Complete descendant \diffmah{} summarizes how a
halo reaches $z=0.4$; it is not automatically the most informative coordinate
system for the state of its progenitor at $z=2$. Re-anchoring the same class of
mass-history model at the prediction epoch aligns the mass normalization,
transition time, and growth indices with the target galaxy state.

The improvement also changes how ``portability'' should be described. The
functional form is portable in the sense that any halo with a resolved past
history can be compressed to four \diffmah{} parameters. The fitted
halo-to-galaxy coefficients are not yet externally portable: they have been
trained and tested only within TNG300.

\subsection{Relation to the deposition-and-migration model}

The direct conversion and the deposition model answer related but different
questions. The deposition model proposes an explicit history-to-radius
mechanism and can be interrogated in terms of when stellar mass arrives and
where it moves. Its assumptions are consequently more numerous and its
parameter degeneracies can be physically informative. The direct model asks
for the lowest-complexity conditional relation that preserves observable CoGs.
It has fewer mechanistic assumptions but does not identify a stellar mass flux
or migration trajectory.

The two models can therefore act as mutual controls. The direct map supplies a
statistical performance ceiling and reveals which profile coordinates contain
assembly information. The deposition model can test whether an explicit
phenomenological mechanism reproduces those dependencies and their redshift
evolution. Agreement on a dependence would be stronger evidence than either
model alone; disagreement would identify where mechanistic assumptions or
statistical shortcuts matter.

\subsection{Why symbolic regression did not yet deliver the hoped-for equation}

Symbolic regression is most useful when the relation is both low-dimensional
and stable under resampling. At $z=0.4$, three recurring residual corrections
recover a substantial fraction of the dense nonlinear gain. At $z=2$, the
nonlinearity is larger but is distributed across more target-specific terms,
and the $z=0.4$ terms transfer weakly. This pattern can arise if the natural
variables change with epoch, if interactions are genuinely numerous, or if a
small symbolic law is obscured by residual covariance and measurement noise.

The correct conclusion is not that symbolic regression is unsuitable. It is
that the current standardized variables and independent-epoch formulation do
not support a compact universal equation. A future search should follow, not
precede, construction of physically consistent epoch-local coordinates and an
explicit redshift variable.

\subsection{Priority of the next tests}

The next decisive analysis is held-epoch closure for the epoch-local model.
Because $\logten M_p$ and $u_c$ are anchored at different cosmic times, the
features require a common physical scaling before their coefficients can be
interpolated. A sensible sequence is:
\begin{enumerate}
  \item express local times relative to the cosmic age at prediction and masses
  relative to concurrent halo mass where appropriate;
  \item standardize with one training-set convention shared across epochs;
  \item remove each interior epoch and predict it from the others;
  \item test whether explicit low-order redshift terms retain the high-$z$
  degree-2 gain;
  \item replace the one-$\rho$ residual model only if a more flexible process
  measurably improves both adjacent and endpoint coherence;
  \item validate the frozen design in another simulation or resolution before
  library graduation.
\end{enumerate}

